\subsection{The exact comparison with coordinate plumbing}

\begin{proposition}[Proved: toric welding is plumbing with the retained coordinate germs]\label{A1-welding-plumbing}
For the corrected surgery data of Proposition~\ref{A1-retained-discs}, every seam is an ordinary analytic plumbing, using the coordinate germs supplied by its two pointed caps. This assertion identifies the welded surface itself. It does not identify the finite plumbing parameter, in a fixed independent choice of germs, with an arbitrary element of the whole normal tensor line.
\end{proposition}
\begin{proof}
Choose centred unit-disc coordinates $u,v$ on the two caps, so the toric boundary identification is $uv=c$, $|c|=1$. Analyticity of the boundaries extends both coordinates to collars of $|u|=|v|=1$. Choose $R>1$ close enough to one that these extensions are embedded and avoid all other caps. In this collar the welding identifies the annuli $R^{-1}<|u|<R$ and $R^{-1}<|v|<R$ by $v=c/u$. The identity map on the remaining component surfaces therefore identifies the welded surface with the plumbing quotient. This works equally for two branches on the same normalization component because their extended cap neighborhoods are disjoint.

More explicitly set $x=a u$, $y=b v$, for any nonzero constants $a,b$. The plumbing equation is $xy=t$ with $t=ab c$. If $z,w$ are reference coordinates at the two centres and $f=u^{-1}$, $g=v^{-1}$, its normal tensor is
\[
 t\,\partial_x\otimes\partial_y
 =ab c\,\frac{f'(0)}a\frac{g'(0)}b
            \partial_z\otimes\partial_w
 =c f'(0)g'(0)\,\partial_z\otimes\partial_w.
\]
This is exactly the tensor of the toric surgery. Shrinking or enlarging the overlapping annuli leaves the quotient unchanged: extend $u$ to the opposite chart by $c/v$ and identify the overlaps using that same holomorphic function. These maps agree where both are defined, so they glue to inverse biholomorphisms of the two quotients. This direct comparison is also the argument in Kra, Appendix III, \S\S14.2--14.5 \cite{AKra1990}.

The coordinate germs in this construction are part of the recovered cap system and vary with its lengths. Proposition~\ref{A2-schwarzian} shows already in a pant that their second jets vary when the other lengths change. Thus retaining these germs, choosing fixed horocyclic germs, and retaining only first-order normal tensors are three distinct operations. The displayed tensor identity describes their first-order comparison, not an extension of a fixed plumbing chart to all tensors.
\end{proof}

\begin{proposition}[Known with derivation: the scope of Kra's coordinates]\label{A1-kra-comparison}
Fix a maximal pants graph of type $(g,n)$ and put $d=3g-3+n>0$. Kra's horocyclic plumbing parameters range in a domain $D$ satisfying
\[
 \mathbb D_{e^{-2\pi}}^{d}\subset D
                   \subset\mathbb D_{e^{-\pi/2}}^{d}.
\]
They cover the moduli of surfaces of that type, generally many times. On marked Teichm\"uller space the logarithmic coordinates $\tau_j$ are global holomorphic coordinates on their image; the corresponding Dehn twist sends $\tau_j$ to $\tau_j+2$. After quotienting these twists, $t_j=e^{\pi i\tau_j}$ are the plumbing coordinates on $D\cap(\C^\times)^d$. They are not coordinates on an unrestricted normal-tensor bundle.
\end{proposition}
\begin{proof}
These are Theorems~1(a),(b) and~3(a),(b) on pp.~501--503 of \cite{AKra1990}, and the identification following Theorem~3(d). The exponential converts translation by $2$ into identity and converts the stated imaginary-part bounds into the displayed polydisc bounds. Zero parameters describe nodes. The bounded domain and the residual moduli identifications exclude the proposed unrestricted global inverse. Proposition~\ref{A1-welding-plumbing} specifies the comparison when the same full coordinate germs are used.
\end{proof}

\begin{proposition}[Known with derivation: Wolpert's metric comparison]\label{A1-wolpert-comparison}
The local family $xy=t$, $0<|t|<1$, is the annulus $|t|<|x|<1$. Its complete curvature-minus-one metric and core length are
\[
 ds_t^2=\left[
 \frac{\pi}{|\log|t||}
 \csc\!\left(\frac{\pi\log|x|}{\log|t|}\right)
 \left|\frac{dx}{x}\right|\right]^2,
 \qquad L_t=\frac{2\pi^2}{-\log|t|}.
\]
On compact subsets of either punctured limiting branch the metric tends to the cusp metric. For a degenerating family of compact stable curves, Wolpert's Expansion~4.2 compares the global hyperbolic metric with the interpolated cusp/annulus metric; its constants are local in the normalization parameters and coordinate germs. This supplies a metric comparison for plumbing, not a global inverse on normal tensors.
\end{proposition}
\begin{proof}
The formula is \cite[\S0.1, p.~417]{AWolpert1990}. Writing $x=e^{\xi+i\theta}$ gives the hyperbolic strip metric; reflection about $\xi=\tfrac12\log|t|$ fixes its core, and integrating there for $0\le\theta\le2\pi$ gives $L_t$. The expansion $a/\sin a\to1$ gives the cusp limit. Expansion~4.2 and Example~4.3, pp.~445--447, apply after fixing small deformation parameters and plumbing germs. The normalization factors $F'_j(0),G'_j(0)$ rescale the parameters in that source, exactly as in Proposition~\ref{A1-first-order}. Its prescribed-curvature correction is essential for the global surface metric. The complete annulus formula alone does not assert equality with that global metric. The compact-family hypothesis is retained here.
\end{proof}

\begin{proposition}[Known with an exact counterexample: arbitrary plumbing need not be a chart]\label{A1-hinich-counterexample}
Even for $(g,n)=(0,4)$, disjoint coordinate discs and the equation $xy=t$ do not ensure that the plumbing map is locally invertible at every nonzero parameter. An instance of Hinich's counterexample is obtained by
\[
 R=\tfrac32,\quad B=\tfrac34,\quad C=10,\quad
 A=\tfrac{420}{569},\quad \varepsilon=\tfrac1{10},
 \qquad \phi(z)=(z-R)^2.
\]
The critical plumbing parameter is $t=9/11$.
\end{proposition}
\begin{proof}
This is an exact specialization of \cite[\S\S2.1--2.3]{AHinich2010}; we give a direct cross-ratio verification. Let $P_1=\mathbf P^1_z$ have markings $A,B$ and node branch $\infty$, with coordinate $y=9/(10z)$ at that branch. Let $P_2=\mathbf P^1_w$ have markings $c=(C-R)^2=289/4,\infty$ and node branch $R^2=9/4$. On its node disc use
\[
 x=\frac{10}{11}(R-\sqrt w),\qquad
 w=(R-\tfrac{11}{10}x)^2,\quad |x|<1,
\]
with the branch $\sqrt{R^2}=R$. This disc is embedded: the map $\phi$ is injective on $|z|<11/10$, since equal values require $z=w$ or $z+w=3$. The latter is impossible there. Its image avoids $c,\infty$ because both preimages $C,3-C$ lie outside that disc. The other node disc $|y|<1$ avoids $A,B$, since both have modulus less than $9/10$. Thus all required coordinate discs and markings satisfy the plumbing hypotheses.

For $0<|t|<1$ plumb by $xy=t$. Set $k=11t/9$. The global coordinate on the resulting sphere is $w$ on the second piece and $(R-kz)^2$ on the first: these agree on the annular overlap. The first retained piece is $|z|<(9/10)/|t|$, hence $|kz|<11/10$ and its map is injective. Its image is $\phi(\mathbb D_{11/10})$. The second retained image is the complement of $\phi(\overline{\mathbb D}_{(11/10)|t|})$. Their intersection is exactly $\phi(\{(11/10)|t|<|u|<11/10\})$, the plumbing overlap. Thus their images cover $\mathbf P^1$ and they glue to a biholomorphism. The four markings are
\[
 a(t)=(R-kA)^2,\quad b(t)=(R-kB)^2,\quad c=289/4,\quad\infty.
\]
They stay distinct for $0<|t|<1$. A moduli coordinate is
\[
 \lambda(t)=\frac{a(t)-c}{b(t)-c},\qquad
 \lambda'(9/11)=0,\qquad
 \lambda''(9/11)=\frac{271040}{371353867}\ne0.
\]
Consequently the plumbing map ramifies at $9/11$. All rational inequalities, derivatives and the equivalent residue cancellation are checked in \path{code/part_a_plumbing_checks.py}. Our explicit branch assignment also specifies which side receives each cap; it avoids relying on the reversed inside/outside cap wording in the arXiv v2 paragraph on p.~4. Hinich's Corollary~1.1 retains local coordinates sufficiently near the nodal origin and does not contradict this example.
\end{proof}

For Hubbard--Koch the comparison is the universal marked family of \cite[Theorems 9.11 and 10.1, and \S13]{AHubbardKoch2014}: a sufficiently small coordinate plumbing is its pullback by a local analytic classifying map. At a node the map on the transverse tangent lines is the tensor identification in Proposition~\ref{A1-first-order}; using the same component deformations supplies the remaining tangent directions. Finite automorphism quotients must still be taken. Theorem~9.11 is explicitly local, and its following remark excludes global injectivity. Thus neither that theorem nor the other opened plumbing sources repairs the unrestricted target of L2.
